\documentclass[11pt]{article}
\usepackage[margin=1in]{geometry}
\usepackage{hyperref}
\usepackage{parskip}

\hypersetup{
    colorlinks=true,
    linkcolor=blue,
    urlcolor=blue
}

\pagestyle{empty}

\begin{document}

{\Large \textbf{NEPR 208, Introduction to Computational Neuroscience}} \\
{\Large \textbf{1\textsuperscript{st} Year Neuroscience Core, 2026}}

April 20th -- May 8th, M W F 1:30 PM -- 3:20 PM

Shaul Druckmann, shauld@stanford.edu \\
Stephen Baccus, baccus@stanford.edu

TA: Jiaming Lu, lujm@stanford.edu

This module will introduce students to computational and theoretical approaches in neuroscience. Emphasis will be on specific questions and how those questions can be answered with computational methods.

Monday and Wednesday classes will be lectures, on Friday students will work on and discuss problem sets.

Website: \url{https://druckmann-lab.github.io/nepr208}

\subsection*{Week 1, April 20 -- 24}

April 20. Introduction and the Perceptron model (Druckmann) \\
April 22. Neural Encoding (Druckmann) \\
April 24. Work on Problem set 1 in class.

\subsection*{Week 2, April 27 -- May 1}

April 27. Neural Population Analysis 1 (Druckmann) \\
April 29. Neural Population Analysis 2 (Druckmann) \\
May 1. Work on Problem set 2 in class.

\subsection*{Week 3, May 4 -- May 8}

May 4. Adaptation and synaptic plasticity (Baccus) \\
May 6. The Hopfield model of context dependent memory (Druckmann) \\
May 8. Work on Problem set 3 in class.

\end{document}
